\section{The analytic threefold constructed in the candidate}
\label{sec:candidate-geometry}

This section audits the construction itself.  It does not infer existence
from the later topology.  The period family, the three local fillings, and
the gluing are proved in that order.  The notation is the candidate's:
\(B\simeq\mathbb P^1\),
\(B^\circ=B\setminus\{p_0,p_1,p_2\}\), and the local monodromy orders at
\(p_1,p_2\) are \(m_1=3,m_2=4\).

\subsection{The period family}

Let \(\Delta\simeq(\mathbb Z/3)*(\mathbb Z/4)\) be the orientation-preserving
\((3,4,\infty)\)-triangle group acting on a simply connected Riemann surface
\(\mathfrak h_z\), with elliptic generators \(g_1,g_2\) and
\(g_0=(g_1g_2)^{-1}\).  Let \(t\) be the affine coordinate on
\(B\) with values \(0,1,\infty\) at \(p_1,p_2,p_0\), respectively.
Write
\[
 \pi:\mathfrak h_z\longrightarrow B\setminus\{p_0\}
\]
for the orbifold quotient map; its stabilizers at chosen lifts
\(z_j\in\pi^{-1}(p_j)\) are \(\langle g_j\rangle\).  Let
\(\rho_V:\Delta\to\operatorname{GL}(V)\) be the printed representation,
\(\rho_V(g_j)=T_j\), and put
\[
 M_g:=\rho_V(g)^{\mathsf t},\qquad A_g:=M_g^{-1}.
\]
Thus \(M_{gh}=M_hM_g\), \(A_{g_j}=A_j\), and
\(A_{g_0}=M_0\) in the notation of \cref{prop:monodromy-certificate}.

\begin{lemma}[simultaneous elliptic-orientation normalization]
\label{lem:triangle-generator-orientations}
\Verified\ The generators may, and henceforth do, satisfy
\begin{equation}
 g_1'(z_1)=e^{-2\pi i/3},\qquad
 g_2'(z_2)=e^{-2\pi i/4}=-i,
\label{eq:triangle-generator-derivatives}
\end{equation}
while \(g_0=(g_1g_2)^{-1}\) remains parabolic.  More precisely, the Cayley
coordinates
\begin{equation}
 s_j(z)=\frac{z-z_j}{z-\bar z_j}
\label{eq:triangle-cayley-coordinates}
\end{equation}
obey
\begin{equation}
 s_1\circ g_1=e^{-2\pi i/3}s_1,
 \qquad
 s_2\circ g_2=-is_2.
\label{eq:triangle-linearizing-identities}
\end{equation}
The two signs in \cref{eq:triangle-generator-derivatives} are a simultaneous
orientation choice; they are not obtained by shrinking a coordinate disc.
\end{lemma}

\begin{proof}
Realize \(\Delta\) as the orientation-preserving subgroup of the reflection
group of a hyperbolic triangle with counterclockwise vertices \(A,B,C\),
angles \(\pi/3,\pi/4,0\), and opposite sides \(a,b,c\), respectively.  Let
\(\sigma_a,\sigma_b,\sigma_c\) denote reflection in those sides and put
\[
 x=\sigma_b\sigma_c,\qquad
 y=\sigma_c\sigma_a,\qquad
 z=\sigma_a\sigma_b.
\]
The product of reflections in two oriented geodesics meeting at angle
\(\theta\) is the rotation through \(2\theta\).  Hence \(x\) is the
counterclockwise rotation about \(A\) through \(2\pi/3\), \(y\) is the
counterclockwise rotation about \(B\) through \(\pi/2\), and \(z\) is
parabolic at the ideal vertex \(C\); this is the standard triangle-group
construction \cite[Chapters~7 and~10]{Beardon1983}.  Moreover
\[
 xyz=\sigma_b\sigma_c\sigma_c\sigma_a\sigma_a\sigma_b=1.
\]
Set
\[
 z_1=A,\qquad z_2=B,\qquad g_1=x^{-1},\qquad g_2=y^{-1}.
\]
Both finite rotations are thereby reversed, so their complex derivatives at
their fixed points are exactly those in
\cref{eq:triangle-generator-derivatives}.  Since
\(g_2g_1=y^{-1}x^{-1}=z\), one also has
\[
 g_1g_2=g_1zg_1^{-1};
\]
thus \(g_1g_2\), and therefore \(g_0=(g_1g_2)^{-1}\), is parabolic.

This simultaneous reversal is forced once the orientation of \(g_1\) and
the cusp conjugacy class are fixed.  Indeed, in
\(\langle x\rangle*\langle y\rangle\simeq C_3*C_4\) every parabolic element
is conjugate to a nonzero power of the primitive cusp word
\(z=y^{-1}x^{-1}=y^3x^2\).  The two mixed choices would instead give
\[
 (x^{-1}y)^{-1}=y^{-1}x=y^3x,
 \qquad
 (xy^{-1})^{-1}=yx^{-1}=yx^2.
\]
Both words are cyclically reduced of length two.  A nontrivial power
\(z^n\) has cyclically reduced length \(2|n|\); hence a mixed word could be
parabolic only for \(|n|=1\).  But neither \(y^3x\) nor \(yx^2\) is a cyclic
permutation of \(y^3x^2\) or its inverse \(xy\).  Thus neither mixed choice
is parabolic.  The derivative selected at \(z_1\) consequently fixes the
derivative at \(z_2\) as well.

Finally, for an elliptic M\"obius transformation \(g\) fixing
\(z_0\in\mathfrak h_z\), the Cayley map
\(s=(z-z_0)/(z-\bar z_0)\) conjugates \(g\) to a disc automorphism fixing
zero.  Such an automorphism is \(s\mapsto g'(z_0)s\).  Applying this to
\(g_1,g_2\) proves \cref{eq:triangle-linearizing-identities}.
\end{proof}

\begin{proposition}[existence and non-degeneracy of the period matrix]
\label{prop:period-input-audit}
\Verified\ There are holomorphic functions
\[
 \tau:\mathfrak h_z\to\mathfrak h,
 \qquad \mu,\beta:\mathfrak h_z\to\mathbb C
\]
satisfying
\[
\begin{array}{lll}
\tau\circ g_0=\tau-1,&\mu\circ g_0=\mu,&\beta\circ g_0=\beta+1,\\
\tau\circ g_2=-1/\tau,&\mu\circ g_2=1+\mu/\tau,
 &\beta\circ g_2=\beta-3-6\mu^2/\tau,
\end{array}
\tag{2.1}
\]
and the corresponding order-three transformation laws obtained from
\(\tau\circ g_1=(\tau-1)/\tau\).  They may be chosen so that
\[
 \operatorname{Im}\tau>0,
 \qquad
 D:=\operatorname{Im}\beta-
 \frac{6(\operatorname{Im}\mu)^2}{\operatorname{Im}\tau}<0.
\tag{2.2}
\]
Consequently
\[
 \Pi=[Z\mid I_2],\qquad
 Z=\begin{pmatrix}6\mu&\tau\\ \beta&\mu\end{pmatrix}
\tag{2.3}
\]
maps the marked rank-four lattice \(\Lambda\) isomorphically onto a lattice
in \(\mathbb C^2\) at every point.  The quotients
\(\mathbb C^2/\Pi(z)\Lambda\) form a holomorphic family of compact
two-dimensional complex tori over \(\mathfrak h_z\), equivariant for
\(\Delta\).
\end{proposition}

\begin{proof}
Normalize the modular function by
\(j(q)=q^{-1}(1+O(q))\) and put
\(J_B=1728(t\circ\pi)\); the modular facts used here are collected in
\cite[Chapter~VII]{Serre1973}.  Put
\[
 \mathfrak h_z^\dagger=
 \mathfrak h_z\setminus J_B^{-1}\{0,1728\},\qquad
 \mathfrak h_\tau^\dagger=
 \mathfrak h\setminus j^{-1}\{0,1728\}.
\]
Away from the deleted sets, the classical modular map
\(j:\mathfrak h_\tau^\dagger\to\mathbb C\setminus\{0,1728\}\) is an
unramified covering with deck group
\(\operatorname{PSL}_2(\mathbb Z)\).  The deleted set in
\(\mathfrak h_z\) is locally finite.  Since \(\mathfrak h_z\) is simply
connected, \(\pi_1(\mathfrak h_z^\dagger)\) is normally generated by small
positive meridians about its points.  At a translate of \(z_1\), \(J_B\)
has local degree three over \(0\); the corresponding meridian therefore has
deck image equal to the third power of the order-three modular generator.
At a translate of \(z_2\), \(J_B\) has local degree four over \(1728\); its
meridian has deck image equal to the fourth power of the order-two modular
generator.  Both images are the identity.  The covering-space lifting
criterion gives a global lift \(\tau\) with \(j(\tau)=J_B\).

Near a deleted point, compose the lift with a biholomorphism
\(\chi:\mathfrak h\to\mathbb D\).  The bounded map \(\chi\circ\tau\)
extends by the removable-singularity theorem.  Its value at the puncture has
modulus strictly less than one: equality would give an interior maximum and
force \(\chi\circ\tau\) to be constant, although its punctured values lie in
 \(\mathbb D\).  Thus the extension still lands in \(\mathfrak h\).

Because \(J_B\) is \(\Delta\)-invariant, \(\tau\circ g\) and \(\tau\) are
lifts of the same map through the Galois covering \(j\) on the punctured
surface.  Hence there is a unique deck transformation
\(\varphi(g)\in\operatorname{PSL}_2(\mathbb Z)\) such that
\begin{equation}
 \tau\circ g=\varphi(g)\circ\tau.
\label{eq:period-deck-homomorphism}
\end{equation}
The identity extends across the deleted points.  Uniqueness and
\((\tau\circ g)\circ h=\tau\circ(gh)\) give
\(\varphi(gh)=\varphi(g)\varphi(h)\), so \(\varphi\) is a homomorphism.

Normalize \(\tau(z_1)=\rho=e^{\pi i/3}\).  The local cubic expansion of
\(j\) at \(\rho\), together with the local degree three of \(J_B\), gives a
simple zero of \(\tau-\rho\) at \(z_1\).  At \(z_2\), choose
\(\gamma\in\operatorname{PSL}_2(\mathbb Z)\) with
\(\tau(z_2)=\gamma(i)\).  The local quadratic expansion of
\(j\circ\gamma\) at \(i\), together with the local degree four of \(J_B\),
 gives a double zero of \(\gamma^{-1}\tau-i\) at \(z_2\).  At \(z_1\),
 differentiating \cref{eq:period-deck-homomorphism} and using the simple
 zero gives
 \[
  \varphi(g_1)'(\rho)=g_1'(z_1)=e^{-2\pi i/3}.
 \]
 The stabilizer of \(\rho\) is generated by
 \(S_1(\tau)=(\tau-1)/\tau\), whose derivative at \(\rho\) is
 \(e^{-2\pi i/3}\); hence \(\varphi(g_1)=S_1\), rather than \(S_1^{-1}\).
 The order-four generator cannot have trivial modular image:
otherwise the invariant local power series \(\gamma^{-1}\tau-i\) would have
order divisible by four, contrary to its order two.  Its modular image
 therefore has order two.  If \(P=\varphi(g_0)\), the exact relation
 \(g_1g_2g_0=1\) gives
 \begin{equation}
  \varphi(g_2)=S_1^{-1}P^{-1}.
 \label{eq:period-second-generator-from-cusp}
 \end{equation}

It remains to determine the cusp image.  Choose a punctured cusp disc
\(D_0^*\subset B\) and a component
\[
 U_0\subset\pi^{-1}(D_0^*)
\]
stabilized by \(\langle g_0\rangle\).  Choose \(R\) and shrink \(D_0^*\) so
that \(\tau(U_0)\) lies in a single component
\[
 K\subset j^{-1}(\{|w|>R\}).
\]
Put \(P=\varphi(g_0)\).  Since \(g_0U_0=U_0\),
\[
 P\tau(U_0)=\tau(g_0U_0)=\tau(U_0).
\]
Hence \(P(K)=K\): the two components \(K\) and \(P(K)\) intersect in the
nonempty set \(\tau(U_0)\).  Let \(c\) be the rational boundary cusp of
\(K\), choose
\(\gamma_0\in\operatorname{PSL}_2(\mathbb Z)\) with
\(\gamma_0(\infty)=c\), and set
\[
 \tau'=\gamma_0^{-1}\tau|_{U_0},\qquad
 P'=\gamma_0^{-1}P\gamma_0.
\]
Then
\[
 \tau'\circ g_0=P'\circ\tau',
\]
and \(P'\) preserves \(K'=\gamma_0^{-1}K\), whose boundary cusp is
\(\infty\).

Suppose \(P'=1\).  Then \(\tau'\) is \(g_0\)-invariant and descends to a
holomorphic map \(\bar\tau':D_0^*\to\mathfrak h\).  Put
\[
 q'=\exp(2\pi i\bar\tau').
\]
After increasing \(R\), the \(q\)-expansion is valid on \(K'\), and
\(0<|q'|<1\).  Thus \(q'\) extends holomorphically across the puncture.
If \(|q'(0)|=1\), the maximum-modulus principle makes \(q'\), hence
\(J_B\), constant.  If \(q'(0)\ne0\), the expansion
\(j(q')=(q')^{-1}(1+O(q'))\) makes \(J_B\) bounded, contrary to its pole.
Therefore \(q'(0)=0\).  But a holomorphic zero has positive winding on a
small positive circle, while \(q'\) has the single-valued holomorphic
logarithm \(2\pi i\bar\tau'\) on \(D_0^*\), so its winding is zero.  This
contradiction proves \(P'\ne1\).  Since \(P'\) preserves the selected cusp
component \(K'\), it is parabolic; hence the original \(P\) is a nontrivial
parabolic.

We now discard the auxiliary conjugation by \(\gamma_0\) and return to the
original normalization
\[
 \tau(z_1)=\rho,\qquad \varphi(g_1)=S_1.
\]
Choose the trace-\(+2\) lift
\(\widetilde P\in\operatorname{SL}_2(\mathbb Z)\) of \(P\).  For a
primitive fixed vector \((p,r)^{\mathsf t}\), it has the unique form
\[
 \widetilde P
 =I+k\begin{pmatrix}-pr&p^2\\-r^2&pr\end{pmatrix},
 \qquad p,r\in\mathbb Z,\quad \gcd(p,r)=1,\quad k\ne0.
\]
With
\[
 \widetilde S_1=\begin{pmatrix}1&-1\\1&0\end{pmatrix},
\]
one has
\[
 \operatorname{tr}(\widetilde S_1^{-1}\widetilde P^{-1})
 =1+k(p^2-pr+r^2).
\]
Since \(S_1^{-1}P^{-1}=\varphi(g_2)\) has projective order two, this trace
is zero.  Hence
\[
 -k(p^2-pr+r^2)=1,
\]
so \(k=-1\) and \(p^2-pr+r^2=1\).  Modulo simultaneous sign, the three
solutions are cycled transitively by \(S_1\).  Conjugating the entire deck
representation by a suitable power of \(S_1\), which fixes \(\rho\) and
leaves \(S_1\) unchanged, therefore gives
\[
 P(w)=w-1,\qquad \varphi(g_2)(w)=-\frac1w.
\]
Rename the resulting period map \(\tau\).  Since \(z_2\) is fixed by the
order-four generator, the new normalization has \(\tau(z_2)=i\), and the
preceding local ramification calculation says exactly that \(\tau-i\) has
order two at \(z_2\).  Let \(t_c\) vanish at \(p_0\) and define anew
\[
 q=e^{2\pi i\tau}.
\]
Since \(\tau\circ g_0=\tau-1\), \(q\) descends to \(D_0^*\).  The simple
pole of \(J_B=j(\tau)\) gives
\[
 q=t_cu_1(t_c),\qquad u_1(0)\ne0.
\]

Writing
\[
 S_1(w)=\frac{w-1}{w},\qquad S(w)=-\frac1w,
 \qquad P(w)=w-1,
\]
the normalized deck certificate is therefore
\begin{equation}
 \varphi(g_1)=S_1,\qquad
 \varphi(g_2)=S,\qquad
 \varphi(g_0)=P,\qquad
 S_1SP=1
 \quad\text{in }\operatorname{PSL}_2(\mathbb Z).
\label{eq:period-deck-certificate}
\end{equation}
Differentiation at \(z_1\), where \(\tau-\rho\) has order one, gives
\(g_1'(z_1)=S_1'(\rho)=\rho^{-2}=e^{-2\pi i/3}\).  At \(z_2\), comparison
of the quadratic leading terms gives only
\[
 \bigl(g_2'(z_2)\bigr)^2=S'(i)=-1.
\]
The choice \(g_2'(z_2)=-i\), rather than \(+i\), is supplied by the
simultaneous parabolic normalization in
\cref{lem:triangle-generator-orientations}; it does not follow from the
quadratic ramification calculation alone.

We now solve the two affine cocycle problems, including the orbifold and
cusp regularity which the abbreviated argument had suppressed.  Define
\(\widetilde\jmath:\Delta\times\mathfrak h_z\to\mathbb C^*\) on the
generators by
\[
 \widetilde\jmath(g_1,z)=-\tau(z),\qquad
 \widetilde\jmath(g_2,z)=\tau(z),
\]
and extend it by
\(\widetilde\jmath(gh,z)=\widetilde\jmath(g,hz)
\widetilde\jmath(h,z)\).  This is well defined: the products around the
two finite cyclic subgroups are
\[
 (-1)^3\tau\frac{\tau-1}{\tau}\frac{-1}{\tau-1}=1,
 \qquad
 \tau\left(-\frac1\tau\right)\tau\left(-\frac1\tau\right)=1.
\]
Moreover \(\widetilde\jmath(g_0,\cdot)=1\).

For an open \(U\subset B\), set
\(\widetilde U=\pi^{-1}(U\setminus\{p_0\})\).  Define
\[
 \mathcal L(U)=\left\{\nu\in\mathcal O(\widetilde U):
 \nu\circ g=\nu/\widetilde\jmath(g,\cdot)\quad(g\in\Delta)\right\}
\]
and let \(\mathcal M(U)\) be the holomorphic functions on
\(\widetilde U\) satisfying
\[
 \mu\circ g_1=\frac{1-\mu}{\tau},\qquad
 \mu\circ g_2=1+\frac\mu\tau.
 \tag{2.1a}
\]
If \(p_0\in U\), both definitions additionally require that the restriction
to the distinguished \(\langle g_0\rangle\)-invariant cusp component,
viewed as a function of \(t_c\), extend holomorphically across \(t_c=0\).
Because \(\widetilde\jmath(g_0)=1\) and \(\mu\circ g_0=\mu\), this
requirement is meaningful and is stable under restriction.  Thus
\(\mathcal L\) is a line sheaf and \(\mathcal M\), if locally nonempty, is
an affine \(\mathcal L\)-torsor.

Local nonemptiness is explicit.  Over an ordinary small disc, prescribe
\(\mu=0\) on one sheet and extend using (2.1a).  At the cusp prescribe
\(\mu=0\) on the distinguished component.  At the elliptic points use
\[
 \mu_1=\frac{2-\tau}{3}\quad\hbox{on }\Delta_1,
 \qquad
 \mu_2=\frac{1-\tau}{2}\quad\hbox{on }\Delta_2.
 \tag{2.1b}
\]
Indeed
\[
 \mu_1\circ g_1=\frac{\tau+1}{3\tau}
 =\frac{1-\mu_1}{\tau},
 \qquad
 \mu_2\circ g_2=\frac{\tau+1}{2\tau}
 =1+\frac{\mu_2}{\tau}.
\]
The cyclic products above, or direct iteration of (2.1a), show that these
prescriptions close after three and four steps.

To identify \(\mathcal L\), use that \(E_6\) has a simple zero at the
order-two elliptic orbit and no other zero in \(\mathfrak h\).  Since
\(\tau-i\) has order two at every translate of \(z_2\), the function
\(E_6\circ\tau\) has only even-order zeros on the simply connected surface
\(\mathfrak h_z\).  Fix a holomorphic square root \(h\),
\(h^2=E_6\circ\tau\), and put
\[
 \Theta=\frac{(E_4\circ\tau)^2h}{\Delta_{\rm mod}\circ\tau}.
\]
For \(j=1,2\), the quotient
\(h(g_jz)/(\tau(z)^3h(z))\) has square one away from its removable zero
set, hence is a constant sign.  At \(z_1\), \(h(z_1)\ne0\) and
\(\rho^3=-1\), so the sign for \(g_1\) is \(-1\).  At \(z_2\), choose the
linearizing coordinate with \(s_2\circ g_2=-is_2\).  The function \(h\) has
a simple zero there, and therefore
\(h(g_2z)/h(z)\to-i=i^3\); the sign for \(g_2\) is \(+1\).  The modular
transformation formulae
\[
 E_4(\gamma\tau)=(c\tau+d)^4E_4(\tau),\quad
 E_6(\gamma\tau)=(c\tau+d)^6E_6(\tau),\quad
 \Delta_{\rm mod}(\gamma\tau)=(c\tau+d)^{12}\Delta_{\rm mod}(\tau)
\]
now give the exact identities
\[
 \Theta\circ g_1=-\Theta/\tau,\qquad
 \Theta\circ g_2=\Theta/\tau.
 \tag{2.1e}
\]
These are precisely the automorphy laws defining \(\mathcal L\).  Since
\(E_4\) has a simple zero at \(\rho\), whereas \(E_6\) has a simple zero
at \(i\), the orders of \(\Theta\) at \(z_1,z_2\) are respectively two
and one.  At the distinguished cusp, \(E_4,E_6\) are units and
\(\Delta_{\rm mod}=q\cdot(\text{unit})\), with
\(q=t_c\cdot(\text{unit})\), so
\(\Theta=t_c^{-1}\cdot(\text{unit})\).
Taking the stabilizer orders into account, division by \(\Theta\) identifies
the ordinary and elliptic stalks with \(\mathcal O_B\), while the cusp stalk
is \(\mathfrak m_{p_0}\Theta\).  Hence
\[
 \mathcal L\simeq\mathcal O_B(-p_0)\simeq\mathcal O_{\mathbb P^1}(-1).
\]
Consequently \(H^1(B,\mathcal L)=H^0(B,\mathcal L)=0\): the locally
nonempty torsor \(\mathcal M\) has a unique global section.  Evaluating at
the fixed points gives
\[
 \mu(z_1)=\frac{2-\rho}{3},\qquad
 \mu(z_2)=\frac{1-i}{2},
\]
and the cusp condition says precisely that its invariant branch extends
holomorphically.

For the second affine problem put
\[
 \phi_1=2-\frac{6(1-\mu)^2}{\tau},\qquad
 \phi_2=-3-\frac{6\mu^2}{\tau}.
\]
Substitution around the finite orbits gives
\[
 \sum_{k=0}^{2}\phi_1\circ g_1^k=0,
 \qquad
 \sum_{k=0}^{3}\phi_2\circ g_2^k=0.
 \tag{2.1c}
\]
For example, the three summands in the first identity are
\[
2-\frac{6(1-\mu)^2}{\tau},\quad
2-\frac{6(\tau-1+\mu)^2}{\tau(\tau-1)},\quad
2+\frac{6\mu^2}{\tau-1},
\]
 Putting them over the common denominator \(\tau(\tau-1)\), the three
 constant terms contribute \(6\tau(\tau-1)\), while the remaining terms
 contribute \(-6\tau(\tau-1)\); hence the total numerator is zero.  For the
 order-four identity the
bracket after extracting \(6/\tau\) is \(2\tau\), cancelling \(-12\).

Define \(\mathcal B(U)\) to consist of holomorphic \(\beta\) on
\(\widetilde U\) with
\(\beta\circ g_j-\beta=\phi_j\); at the cusp require
\(b=\beta+\tau\) to extend holomorphically in \(t_c\).  Differences of
sections are ordinary holomorphic functions on \(U\), so
\(\mathcal B\) is a pseudo-torsor under \(\mathcal O_B\).  It is locally
nonempty: prescribe \(\beta=0\) on one ordinary sheet; prescribe
\(\beta=-\tau\) on the distinguished cusp component; and at \(p_j\) use
\[
 \beta_j=\frac1{m_j}\sum_{k=0}^{m_j-1}
 k\,(\phi_j\circ g_j^k).
 \tag{2.1d}
\]
Equation (2.1c) makes the last expression telescope to
\(\beta_j\circ g_j-\beta_j=\phi_j\).  Therefore \(\mathcal B\) is an
\(\mathcal O_B\)-torsor.  Since \(H^1(B,\mathcal O_B)=0\), it has global
sections, and any two differ by a constant; by construction
\(b=\beta+\tau\) is invariant and holomorphic at the cusp.

The real determinant is
\[
 \det_{\mathbb R}\Pi=-\det\operatorname{Im}Z
 =\operatorname{Im}\tau\,D.
\tag{2.4}
\]
For completeness, the equivariance can be checked before any quotient is
taken.  Regard the rows of \(\Pi\) as the two covectors
\[
 \sigma_1=6\mu\gamma+\tau u+w,
 \qquad
 \sigma_2=\beta\gamma+\mu u+\delta
 \quad\text{in }V\otimes\mathbb C.
 \tag{2.3a}
\]
Multiplication by the printed matrices gives
\[
\begin{aligned}
T_1\sigma_1&=(6\mu-6)\gamma+(1-\tau)u-\tau w,\\
T_1\sigma_2&=(\beta+2)\gamma+(1-\mu)u+(1-\mu)w+\delta,\\
T_2\sigma_1&=(6\mu+6\tau)\gamma-u+\tau w,\\
T_2\sigma_2&=(\beta+6\mu-3)\gamma+u+\mu w+\delta.
\end{aligned}
\tag{2.3b}
\]
Using the four affine transformation laws in (2.1a)--(2.1d), coefficient
comparison in (2.3b) yields
\[
\begin{aligned}
 \sigma_1\circ g_1&=-\tau^{-1}T_1\sigma_1,
 &\sigma_2\circ g_1&=T_1\sigma_2+(1-\mu)\tau^{-1}T_1\sigma_1,\\
 \sigma_1\circ g_2&=\tau^{-1}T_2\sigma_1,
 &\sigma_2\circ g_2&=T_2\sigma_2-\mu\tau^{-1}T_2\sigma_1.
\end{aligned}
\tag{2.3c}
\]
For instance, dividing the first line of (2.3b) by \(-\tau\) gives the
coefficients \(6(1-\mu)/\tau\), \((\tau-1)/\tau\), and \(1\), exactly
those of \(\sigma_1(g_1z)\); subtracting \((1-\mu)\) times that normalized
row from \(T_1\sigma_2\) gives \(\sigma_2(g_1z)\).  The two \(g_2\)
identities follow by the same coefficient check.  Thus, with
\[
R_{g_1}=\begin{pmatrix}-1/\tau&0\\(1-\mu)/\tau&1\end{pmatrix},
\qquad
R_{g_2}=\begin{pmatrix}1/\tau&0\\-\mu/\tau&1\end{pmatrix},
\]
pairing (2.3c) with \(\lambda\in\Lambda\) gives
\[
 \Pi(gz)=R_g(z)\Pi(z)M_g.
 \tag{2.3d}
\]
Since \(M_{gh}=M_hM_g\), uniqueness of the left factor gives
\(R_{gh}(z)=R_g(hz)R_h(z)\); hence these formulae define an action, not
merely generator-wise maps.  Taking real determinants in (2.3d), using
\(\det M_g=1\) and
\(\operatorname{Im}\tau(g_jz)=\operatorname{Im}\tau/|\tau|^2\), shows
\(D(gz)=D(z)\).  At the cusp,
\[
 D=\operatorname{Im}b-\operatorname{Im}\tau-
 \frac{6(\operatorname{Im}\mu)^2}{\operatorname{Im}\tau}
 \longrightarrow-\infty.
\]
The descended function is bounded above on the complement of a cusp
neighbourhood.  Replacing \(\beta\) by \(\beta+c_0\) adds
\(\operatorname{Im}c_0\) to \(D\), so a sufficiently negative imaginary
constant makes (2.2) hold globally.  Equation (2.4) then makes \(\Pi\) a real
isomorphism everywhere.  The equivariance identity supplies the claimed
\(\Delta\)-action and completes the construction of the period family.
\end{proof}

For the cusp filling, record the local periods explicitly.  Let \(t_c\) be
the coordinate vanishing at \(p_0\).  The simple pole of \(j\) gives
\[
 q=e^{2\pi i\tau}=t_cu_1(t_c),\qquad u_1(0)\ne0.
\]
After choosing a logarithm of \(u_1\), put
\[
 h=(2\pi i)^{-1}\log u_1,\qquad s=\tau-h,
 \qquad b=\beta+\tau.
\]
Then \(e^{2\pi is}=t_c\), \(s\circ g_0=s-1\), and \(\mu,b,h\) are
holomorphic functions of \(t_c\) at the origin.  Consequently
\[
 B_0=\begin{pmatrix}0&1\\-1&0\end{pmatrix},\qquad
 Z=sB_0+C(t_c),\qquad
 C(t_c)=
 \begin{pmatrix}
 6\mu(t_c)&h(t_c)\\
 b(t_c)-h(t_c)&\mu(t_c)
 \end{pmatrix},
 \tag{2.5a}
\]
and \(C\) is holomorphic on a full cusp disc.  In particular, for
\(\bar\lambda\in\bar\Lambda\simeq\mathbb Z^2\),
\[
 \exp(2\pi iZ\bar\lambda)
 =t_c^{B_0\bar\lambda}c_{\bar\lambda}(t_c),
 \qquad
 c_{\bar\lambda}(t_c)=\exp(2\pi iC(t_c)\bar\lambda).
 \tag{2.5b}
\]

\subsection{The toric filling at the cusp}

Let \(\mathcal T_{A_2}\) be the translation-invariant triangulation of
\(\mathbb R^2\) by the two triangles
\(\{v,v+e_1,v+e_2\}\) and
\(\{v+e_1,v+e_2,v+e_1+e_2\}\).  In
\(N'=\mathbb Z^2\oplus\mathbb Z\), take the cones over all cells at height
one, and denote the associated toric threefold by \(Y\).
We use the standard orbit--cone and affine-chart constructions of
\cite[Chapters~1 and~3]{Fulton1993}; all nonstandard assertions caused by
the infinite fan are proved directly below.

For \(\bar\lambda\in\bar\Lambda\), define the integral linear map
\[
 \varphi_{\bar\lambda}(y,y_3)
 =(y+y_3B_0\bar\lambda,y_3).
\]
On the height-one slice this is translation by
\(B_0\bar\lambda\in\mathbb Z^2\), so it preserves the triangulation and
maps the cone over a cell \(P\) to the cone over
\(P+B_0\bar\lambda\).  It therefore induces a toric automorphism
\(\Phi_{\bar\lambda}\) of \(Y\), satisfying
\[
 \Phi_{\bar\lambda}(x,t)
 =(t^{B_0\bar\lambda}x,t)
\]
on the dense torus.  Because \(t\) and \(C(t)\) are holomorphic on the
full cusp model, multiplication by
\((c_{\bar\lambda}(t),1)\in(\mathbb C^*)^3\) defines a biholomorphism on
all of \(Y_\epsilon\).  Set
\[
 \Psi_{\bar\lambda}(p)
 =(c_{\bar\lambda}(t(p)),1)\cdot\Phi_{\bar\lambda}(p).
\]
The maps \(\Phi_{\bar\lambda}\) fix the first two cocharacter directions
and preserve \(t\), so they commute with this torus multiplication.  The
linearity of \(C(t)\bar\lambda\) gives
\(c_{\bar\lambda}c_{\bar\mu}=c_{\bar\lambda+\bar\mu}\); hence
\(\Psi_{\bar\lambda}\Psi_{\bar\mu}
=\Psi_{\bar\lambda+\bar\mu}\).  This proves, before taking any quotient,
both the holomorphic extension and the action law used below.

For a triangle \(\sigma=\{v_0,v_1,v_2\}\), let
\(\theta_i^\sigma\) be its barycentric coordinates and let
\(z_i=\chi^{m_i}\) be the coordinates dual to
\((v_i,1)\).  Set
\[
 \sigma^{(\eta)}=\{y:\theta_i^\sigma(y)>-\eta\text{ for }i=0,1,2\},
 \quad
 \mathcal U_\sigma(S)=\{|z_i|<S\text{ for all }i\}.
\]

\begin{lemma}[toric chart control]
\label{lem:toric-chart-control}
For a torus point \(p=(x,t)\) with \(0<|t|<1\), put
\[
 y(p)=\frac{(\log|x_1|,\log|x_2|)}{\log|t|}.
\]
Then
\[
 p\in\mathcal U_\sigma(S)
 \quad\Longleftrightarrow\quad
 y(p)\in\sigma^{(\log S/|\log|t||)}.
 \tag{2.5c}
\]
Every point of \(\{|t|<1\}\subset Y\) belongs to
\(\overline{\mathcal U}_\sigma(1)=\{|z_i|\le1\}\) for some \(\sigma\).
Consequently the sets \(\mathcal U_\sigma(2)\) cover \(Y_\epsilon\), and
every compact subset is contained in the union of finitely many members of
this cover.
\end{lemma}

\begin{proof}
Since \((y,1)=\sum_i\theta_i^\sigma(y)(v_i,1)\), duality gives
\[
 \log|z_i(p)|=\log|t(p)|\,\theta_i^\sigma(y(p)).
\]
Because \(\log|t|<0\), this is exactly (2.5c).

For the boundary assertion let \(p\) lie in the toric orbit
\(O_{\widehat P}\) of the cone over a cell \(P\).  In a chart with
\(\sigma\supseteq P\), the coordinates corresponding to vertices of
\(P\) vanish.  On the quotient lattice
\(N'/(N'\cap\operatorname{span}\widehat P)\), the remaining coordinates
are the characters dual to the image cone \(\bar\sigma\).  Hence
\(p\in\overline{\mathcal U}_\sigma(1)\) exactly when
\(-\operatorname{Log}_{\widehat P}(p)\in\bar\sigma\).
The cones \(\bar\sigma\), for triangles \(\sigma\supseteq P\), cover the
quotient real vector space: lift a vector \(n\), choose a relative interior
point \(p_0\in P\times\{1\}\), and for sufficiently small \(u>0\) place
the height-one rescaling of \(p_0+un\) in a triangle of the locally finite
triangulation.  That triangle contains \(P\), and its quotient cone contains
\(n\).  This proves the asserted cover.  Compactness then supplies a finite
subcover.
\end{proof}

\begin{theorem}[the cusp quotient and its central fibre]
\label{thm:candidate-cusp-geometry}
\Verified\ For the period functions and holomorphic cusp matrix \(C(t)\)
constructed above, and sufficiently small \(\epsilon>0\), the maps
\[
 \Psi_{\bar\lambda}(x,t)=
 \bigl(c_{\bar\lambda}(t)t^{B_0\bar\lambda}x,t\bigr),
 \qquad
 c_{\bar\lambda}(t)=\exp(2\pi iC(t)\bar\lambda),
\tag{2.5}
\]
extend to a free properly discontinuous action of
\(\bar\Lambda\simeq\mathbb Z^2\) on
\(Y_\epsilon=\{|t|<\epsilon\}\).  The quotient
\[
 f_0:N_0=Y_\epsilon/\bar\Lambda\longrightarrow\Delta_\epsilon
\]
is a proper holomorphic map from a connected smooth complex threefold.
Its nonzero fibres are the tori of \cref{prop:period-input-audit}.  Its full
scheme-theoretic central fibre \(W\) is reduced and has local equations
\[
 z_0,\qquad z_0z_1,\qquad z_0z_1z_2.
\tag{2.6}
\]
The surface \(W\) is irreducible and nonnormal.  Its normalization is
\(\nu:E_0\simeq\mathrm{dP}_6\to W\), obtained by identifying opposite sides
of the toric hexagon; its double locus consists of three rational curves
through exactly two triple points.  Thus \(W\) is a normal-crossings surface,
but it is not a simple-normal-crossings divisor.  Moreover
\[
 \omega_{N_0}\simeq\mathcal O_{N_0}.
\tag{2.6a}
\]
\end{theorem}

\begin{proof}
For a bounded cell \(P\), the cone over \(P\times\{1\}\) is strongly
convex, and two such cones meet in the cone over the common face of the two
cells.  The locally finite triangulation therefore defines a countable
rational fan with support
\((\mathbb R^2\times\mathbb R_{>0})\cup\{0\}\).  For the determinant
calculation, order the height-one generators of the two kinds as
\[
 ((v,1),(v+e_1,1),(v+e_2,1)),\qquad
 ((v+e_1,1),(v+e_2,1),(v+e_1+e_2,1)).
\]
Their determinants are respectively \(+1\) and \(-1\), hence have absolute
value one.  Every cone is therefore unimodular.  The ordinary pairwise affine-chart
separation proof for a toric fan applies without a finiteness hypothesis;
countability of the fan gives second countability.  Thus \(Y\) is a
connected Hausdorff smooth complex threefold with dense torus.

If \(m_0,m_1,m_2\) is the basis dual to the height-one generators, then the
corresponding affine chart is \(\mathbb C^3\), with coordinates
\(z_i=\chi^{m_i}\), and
\[
 t=z_0z_1z_2.
\tag{2.7}
\]
Since every primitive ray generator has height one, the divisor of \(t\) is
\(\sum_vE_v\), all with multiplicity one.  Equation (2.7) proves that the
scheme-theoretic fibre \(t^{-1}(0)\) is reduced and normal crossings.
The star of a height-one vertex has rays
\(e_1,e_2,e_2-e_1,-e_1,-e_2,e_1-e_2\), the complete smooth hexagonal fan;
hence every irreducible component \(E_v\) is \(\mathrm{dP}_6\).

The local canonical trivialization is also intrinsic to this toric model.
On the dense torus put
\[
 \Omega=\frac{dx_1\wedge dx_2\wedge dt}{x_1x_2}.
\]
If \((z_0,z_1,z_2)\) is a unimodular toric chart, then
\[
 \bigwedge_i\frac{dz_i}{z_i}
 =\pm\frac{dx_1}{x_1}\wedge\frac{dx_2}{x_2}\wedge\frac{dt}{t},
 \qquad t=z_0z_1z_2,
\]
and therefore \(\Omega=\pm dz_0\wedge dz_1\wedge dz_2\).  It extends
holomorphically and nowhere vanishingly across every boundary chart.
For a deck transformation \(\Psi_{\bar\lambda}\), formula (2.5) gives
\(\Psi_{\bar\lambda}^*(dx_i/x_i)=dx_i/x_i+\alpha_i(t)dt\) and
\(\Psi_{\bar\lambda}^*dt=dt\); every additional term in the pullback of
\(\Omega\) contains \(dt\) twice.  Hence \(\Omega\) is invariant and descends
to a nowhere-vanishing three-form on \(N_0\), proving the intrinsic
trivialization (2.6a).  No global threefold has been invoked at this stage.

Write \(R(t)=-2\pi\operatorname{Im}C(t)\), which is bounded because \(C\)
is holomorphic at zero.  Choose \(\epsilon\) such that
\(|\log\epsilon|>2\|R\|_\infty\).  At a torus point fixed by a nonzero
\(\bar\lambda\), taking logarithmic absolute values in (2.5) gives
\[
 \log|t|B_0\bar\lambda+R(t)\bar\lambda=0,
\]
impossible because \(|B_0\bar\lambda|=|\bar\lambda|\) and the first term
has more than twice the norm of the second.  On \(t=0\), a cell is translated
by \(B_0\bar\lambda\ne0\); a bounded cell cannot be fixed by a nonzero
translation.  Thus the action is free.

For proper discontinuity put
\[
 B_t=B_0+\frac{R(t)}{\log|t|}.
\]
Then \(\|B_t-B_0\|\le1/2\), so
\(|B_t\bar\lambda|\ge|\bar\lambda|/2\).  In logarithmic position
coordinates, (2.5) is
\[
 y(\Psi_{\bar\lambda}p)=y(p)+B_t\bar\lambda.
\tag{2.8}
\]
Put \(\eta_0=\log2/|\log\epsilon|\).  If
\(\Psi_{\bar\lambda}\mathcal U_\sigma(2)\) meets
\(\mathcal U_{\sigma'}(2)\), the intersection is open and hence contains a
point of the dense invariant torus.  By \cref{lem:toric-chart-control}, for
that point \(p\) and its image \(p'\),
\[
 y(p)\in\sigma^{(\eta_0)},\qquad
 y(p')\in{\sigma'}^{(\eta_0)}.
\]
Both thickenings are bounded, so (2.8) gives
\[
 \frac12|\bar\lambda|
 \le |B_t\bar\lambda|
 =|y(p')-y(p)|
 \le \sup_{a\in{\sigma'}^{(\eta_0)},\,b\in\sigma^{(\eta_0)}}|a-b|.
\]
Thus only finitely many \(\bar\lambda\) can make this fixed pair of charts
meet.  A compact set is contained in finitely many of the charts from
\cref{lem:toric-chart-control}; taking the finite union over all chart pairs
proves proper discontinuity.  Freeness and proper discontinuity give a
Hausdorff complex-manifold quotient and a holomorphic covering map.

To prove properness, fix \(\epsilon_1<\epsilon\).  Let \(\mathcal F_0\) be
the finite set of triangles meeting the closed Euclidean ball of radius
\(\sqrt2\), and put
\[
 K_0=\bigcup_{\sigma\in\mathcal F_0}
 \left(\overline{\mathcal U}_\sigma(1)
 \cap\{|t|\le\epsilon_1\}\right).
\]
This is compact.  For a torus point \(p\) with
\(0<|t(p)|\le\epsilon_1\), put \(u=B_{t(p)}^{-1}y(p)\) and choose
\(\bar\lambda\in\mathbb Z^2\) with
\(|u+\bar\lambda|_\infty\le1/2\).  Since
\(\|B_t\|\le3/2\), equation (2.8) gives
\[
 |y(\Psi_{\bar\lambda}p)|
 =|B_t(u+\bar\lambda)|
 \le\frac{3}{2\sqrt2}<\sqrt2.
\]
Choose a triangle containing this point.  Formula (2.5c) with \(S=1\)
places the translate in the corresponding
\(\overline{\mathcal U}_\sigma(1)\subset K_0\).

If \(p\) lies in the central orbit belonging to a cell \(P\), translate a
vertex of \(P\) to zero; this is possible because
\(B_0\bar\Lambda=\mathbb Z^2\).  The refined boundary part of
\cref{lem:toric-chart-control} then places the translate in
\(\overline{\mathcal U}_\sigma(1)\) for a triangle \(\sigma\) containing
zero, hence in \(K_0\).  Therefore the image of \(K_0\) in the quotient is
exactly \(f_0^{-1}(\{|t|\le\epsilon_1\})\), which is compact.  Every compact
subset of the base lies in such a closed subdisc, so \(f_0\) is proper.

For \(t\ne0\), choose \(s\) with \(e^{2\pi is}=t\).  The multiplier in
(2.5) is \(\exp(2\pi iZ(s)\bar\lambda)\), so exponentiation identifies the
fibre with
\[
 \mathbb C^2/(\mathbb Z^2+Z(s)\mathbb Z^2).
\]
It is compact because \(\operatorname{Im}Z(s)\) is invertible by (2.2).

The quotient map is a local biholomorphism; therefore the local equation
(2.7), its multiplicity one, and reducedness descend to \(W\).  Every cell
has a vertex, and \(B_0\bar\Lambda=\mathbb Z^2\) acts simply transitively on
vertices, so every orbit in the central fibre meets \(E_0\).  Thus
\(E_0\to W\) is surjective.  Two interior points cannot be identified; on
the boundary, precisely opposite edges and the corresponding hexagon
vertices are identified.  The fibres have at most three points.  Since
\(E_0\) is compact, the map is finite; since it is generically injective
from a smooth irreducible surface to the reduced surface \(W\), it is the
normalization.  The two translation-orbits of triangles give the two triple
points, and the three unoriented edge directions give the three double
curves.  Irreducibility follows from the finite surjection by \(E_0\).
\end{proof}

\begin{proposition}[the cusp punctured-overlap biholomorphism]
\label{prop:candidate-cusp-overlap}
\Verified\ Let \(D_0^*=\{0<|t_c|<\epsilon\}\), and shrink it if necessary
so that its inverse image has a distinguished component \(U_0\) stabilized
exactly by \(\langle g_0\rangle\).  Exponentiation induces a
fibre-preserving biholomorphism
\[
 J|_{D_0^*}\xrightarrow{\ \sim\ }
 N_0|_{D_0^*}=N_0\setminus W.
\tag{2.8i}
\]
On the covering space it is given by
\[
 \mathcal E:U_0\times\mathbb C^2\longrightarrow
             (\mathbb C^*)^2\times D_0^*,
 \qquad
 \mathcal E(z,\zeta_1,\zeta_2)
 =\bigl(e^{2\pi i\zeta_1},e^{2\pi i\zeta_2},e^{2\pi is(z)}\bigr).
\tag{2.8ii}
\]
\end{proposition}

\begin{proof}
First identify every deck transformation.  From
\(s\circ g_0=s-1\), the equality
\(t_c\circ\pi=e^{2\pi is}\), and the fact that \(t_c\) is a base
coordinate, \(ds\) never vanishes on \(U_0\).  Thus \(s\) is a local
biholomorphism to the logarithmic cover of \(D_0^*\).  The restriction
\(U_0\to D_0^*\) is a connected regular cover with deck group
\(\langle g_0\rangle\simeq\mathbb Z\).  Under the classification of connected
covers of the punctured disc, it corresponds to a subgroup of
\(\pi_1(D_0^*)\simeq\mathbb Z\); because its deck group is infinite, that
subgroup is zero.  Hence \(U_0\) is the universal cover.  The logarithmic
cover is universal as well, and its deck generator acts by
\(s\mapsto s-1\).  Therefore \(s\) is an equivariant biholomorphism onto that
logarithmic cover.

The cusp transformation laws give
\[
 \Pi(g_0z)=[Z(s-1)\mid I_2]=[Z(s)-B_0\mid I_2].
\]
In the same block decomposition, the explicit matrix in
\cref{prop:monodromy-certificate} gives
\[
 M_{g_0}=M_0^{-1}
 =\begin{pmatrix}I_2&0\\-B_0&I_2\end{pmatrix}.
\]
Consequently
\[
 \Pi(g_0z)=\Pi(z)M_{g_0}.
\]
Comparison with (2.3d) and real invertibility of \(\Pi(g_0z)\) yields
\(R_{g_0}=I_2\).  To make the crossed action explicit, for
\(\lambda\in\Lambda\) and \(g\in\Delta\) put
\[
 T_\lambda(z,\zeta)=(z,\zeta+\Pi(z)\lambda),
 \qquad
 \widetilde g(z,\zeta)=(gz,R_g(z)\zeta).
\]
Equation (2.3d), together with \(A_g=M_g^{-1}\), gives
\[
\begin{aligned}
 \widetilde gT_\lambda(z,\zeta)
 &=\bigl(gz,R_g(z)\zeta+R_g(z)\Pi(z)\lambda\bigr)\\
 &=\bigl(gz,R_g(z)\zeta+\Pi(gz)A_g\lambda\bigr)
  =T_{A_g\lambda}\widetilde g(z,\zeta).
\end{aligned}
\]
Thus
\[
 \widetilde gT_\lambda=T_{A_g\lambda}\widetilde g.
\tag{2.8ii-a}
\]
At the cusp \(A_{g_0}=M_0\), so
\[
 \widetilde g_0(z,\zeta)=(g_0z,\zeta),
 \qquad
 \widetilde g_0T_\lambda=T_{M_0\lambda}\widetilde g_0.
\tag{2.8ii-b}
\]
In particular \(\mathcal E\) is invariant under \(\widetilde g_0\).

Let
\(\Lambda_{\rm tor}=\mathbb Z\widehat w\oplus
\mathbb Z\widehat\delta\).  By \cref{thm:cusp-lattices},
\[
 \Lambda_{\rm tor}=\ker(M_0-I)=\operatorname{im}(M_0-I).
\tag{2.8ii-c}
\]
Put
\[
 \mathcal G_0=\langle T_\lambda,\widetilde g_0:\lambda\in\Lambda\rangle
 \simeq\Lambda\rtimes_{M_0}\langle\widetilde g_0\rangle.
\]
Since \(M_0\) fixes \(\Lambda_{\rm tor}\) pointwise, (2.8ii-b) shows that
\[
 H:=\{T_\nu\widetilde g_0^{\,n}:
          \nu\in\Lambda_{\rm tor},\ n\in\mathbb Z\}
 \simeq\Lambda_{\rm tor}\times\langle\widetilde g_0\rangle.
\tag{2.8ii-d}
\]
Moreover, for \(\lambda\in\Lambda\) and
\(\nu\in\Lambda_{\rm tor}\),
\[
 \widetilde g_0T_\nu\widetilde g_0^{-1}=T_\nu,
 \qquad
 T_\lambda T_\nu T_\lambda^{-1}=T_\nu,
 \qquad
 T_\lambda\widetilde g_0T_\lambda^{-1}
 =T_{(I-M_0)\lambda}\widetilde g_0\in H.
\tag{2.8ii-e}
\]
These generator conjugations prove that \(H\triangleleft\mathcal G_0\).
Translation by \(\Lambda_{\rm tor}\) adds an integral vector to \(\zeta\),
so it lies in the deck group of \(\mathcal E\).
Conversely, if two points have the same image under \(\mathcal E\), their
\(s\)-coordinates differ by an integer and their \(\zeta\)-coordinates by
an element of \(\mathbb Z^2\).  Equivariance and injectivity of the
logarithmic coordinate then show that they differ by a unique element of \(H\).
Thus \(\mathcal E\) identifies
\((U_0\times\mathbb C^2)/H\) with
\((\mathbb C^*)^2\times D_0^*\).

It remains to identify the residual action.  Write
\(\bar\Lambda=\Lambda/\Lambda_{\rm tor}\), and let \(\lambda\) be any lift
of \(\bar\lambda\).  In the period basis, translation by \(\lambda\) adds
\(Z(s)\bar\lambda\) to \(\zeta\), up to an integral vector killed by
exponentiation.  Therefore (2.5b) gives
\[
 \mathcal E(z,\zeta+\Pi(z)\lambda)
 =\bigl(c_{\bar\lambda}(t_c)t_c^{B_0\bar\lambda}x,t_c\bigr)
 =\Psi_{\bar\lambda}\bigl(\mathcal E(z,\zeta)\bigr).
\tag{2.8iii}
\]
This is independent of the chosen lift.  Moreover \(M_0\) is the identity
on \(\bar\Lambda\), by (2.8ii-c).  The crossed multiplication law is
\[
 (T_\lambda\widetilde g_0^{\,n})
 (T_\mu\widetilde g_0^{\,m})
 =T_{\lambda+M_0^n\mu}\widetilde g_0^{\,n+m}.
\]
It follows that
\[
 \vartheta:\mathcal G_0\longrightarrow\bar\Lambda,
 \qquad
 \vartheta(T_\lambda\widetilde g_0^{\,n})=\bar\lambda,
\]
is a surjective homomorphism with kernel \(H\).  Hence
\(\mathcal G_0/H\simeq\bar\Lambda\), and (2.8iii) identifies this residual
action with the displayed \(\Psi_{\bar\lambda}\)-action.  Dividing by it
gives
\[
 (U_0\times\mathbb C^2)/\mathcal G_0
 \simeq
 ((\mathbb C^*)^2\times D_0^*)/\bar\Lambda.
\]
The left side is \(J|_{D_0^*}\), while the right side is
\(N_0|_{D_0^*}\) because \(t^{-1}(\mathbb C^*)\) is the dense torus of
\(Y\).  This proves (2.8i), including holomorphicity of the inverse.
\end{proof}

\subsection{The finite logarithmic fillings and the global gluing}

Set
\[
 \varepsilon=\widehat\gamma+2\widehat u-4\widehat w,
 \qquad
 \varepsilon'=\widehat\gamma+3\widehat u-3\widehat w.
 \tag{2.8a}
\]
Direct multiplication by the matrices in
\cref{prop:monodromy-certificate} gives
\(A_1\varepsilon=\varepsilon\) and
\(A_2\varepsilon'=\varepsilon'\).

A \emph{linearizing pair} \((\Delta_j,s_j)\) consists of a simply connected
\(g_j\)-invariant neighbourhood of \(z_j\), disjoint from all its translates
outside \(\langle g_j\rangle\), and an injective coordinate satisfying
\[
 s_j(z_j)=0,\qquad s_j\circ g_j=\zeta_js_j,
 \qquad \zeta_j=e^{-2\pi i/m_j}.
 \tag{2.8b}
\]
Shrinking gives such a pair.  Put \(D_j=\pi(\Delta_j)\),
\(\Delta_j^*=\Delta_j\setminus\{z_j\}\), and choose the base coordinate
\(t_j\) by \(t_j\circ\pi=s_j^{m_j}\).  Let
\[
 \mathcal T=(\mathfrak h_z\times\mathbb C^2)/\Lambda,
 \qquad
 \mathfrak h_z^\circ=\mathfrak h_z\setminus
 (\Delta z_1\cup\Delta z_2),
 \qquad
 J=(\mathcal T|_{\mathfrak h_z^\circ})/\Delta
 \]
\[
 J'_j=\mathcal T|_{\Delta_j},
 \qquad
 A_0^{(j)}=\mathbb C^2/\Pi(z_j)\Lambda.
 \tag{2.8c}
\]

\begin{lemma}[the middle period quotient is proper]
\label{lem:middle-period-proper}
The induced \(\Delta\)-action on
\(\mathcal T|_{\mathfrak h_z^\circ}\) is free and properly discontinuous.
Consequently \(J\) is a Hausdorff complex manifold and its natural map
\[
 J\longrightarrow B_{\rm mid}:=\mathfrak h_z^\circ/\Delta
\]
is a proper holomorphic submersion with compact torus fibres.
\end{lemma}

\begin{proof}
The deleted translates contain every elliptic fixed point, so the
\(\Delta\)-action on \(\mathfrak h_z^\circ\) is free; it is properly
discontinuous because it is the restriction of the Fuchsian action.  The
action on \(\mathcal T\) projects equivariantly to this base action.  Hence
an element fixing a point of \(\mathcal T|_{\mathfrak h_z^\circ}\) fixes
its base point and is the identity.  If \(K_1,K_2\) are compact subsets of
the total space and \(gK_1\cap K_2\ne\varnothing\), then
\(g\,\operatorname{pr}(K_1)\cap\operatorname{pr}(K_2)\ne\varnothing\).
Base proper discontinuity leaves only finitely many such \(g\), proving
proper discontinuity upstairs and hence the Hausdorff manifold assertion.

For properness, take a simply connected open set \(U\Subset B_{\rm mid}\)
evenly covered by \(\mathfrak h_z^\circ\), and one lift \(\widetilde U\).
The quotient over \(U\) is identified with
\(\mathcal T|_{\widetilde U}\).  Before quotienting, projection is a
holomorphic submersion, and the quotient maps are local biholomorphisms;
hence the displayed map is a holomorphic submersion.  The flat real-analytic trivialization
\[
 \widetilde U\times(\Lambda\otimes\mathbb R)/\Lambda
 \longrightarrow\mathcal T|_{\widetilde U},\qquad
 (z,x+\Lambda)\longmapsto[z,\Pi(z)x]
\]
shows this properness explicitly: if \(C\subset U\) is compact, its chosen
lift \(\widetilde C\subset\widetilde U\) is compact and the inverse image of
\(C\) is homeomorphic to
\(\widetilde C\times(\Lambda\otimes\mathbb R)/\Lambda\), hence compact.
Thus the map is proper over \(U\).  Properness is local on the target, so
the displayed map is proper.
\end{proof}

For \(b\in\Lambda\otimes\mathbb R\), define the varying holomorphic lift
\[
 \widetilde g_j^{b}(z,\zeta)=
 \bigl(g_jz,R_{g_j}(z)\zeta+\Pi(g_jz)b\bigr).
 \tag{2.8d}
\]
Equation (2.3d) says
\[
 R_{g_j}(z)\Pi(z)=\Pi(g_jz)A_j,
 \tag{2.8e}
\]
so (2.8d) normalizes the lattice action and descends to \(J'_j\).  In the
flat real-analytic coordinates
\[
 \Xi(z,x+\Lambda)=[z,\Pi(z)x]
\]
one has, with \(\Sigma_k=I+A_j+\cdots+A_j^{k-1}\),
\[
 \Xi^{-1}(\widetilde g_j^{b})^k\Xi(z,x)
 =\bigl(g_j^kz,A_j^kx+\Sigma_kb\bigr).
 \tag{2.8f}
\]
We call \(\widetilde g_j^{0}\) the linear lift and, for
\(v\in\Lambda^{A_j}\), call
\(\widetilde g_j^{\log}:=\widetilde g_j^{v/m_j}\) the logarithmic lift.

\begin{proposition}[finite fillings]
\label{prop:candidate-finite-fillings}
\Verified\ Let \(A_j\) be the two finite monodromies and choose
\[
 v_1=\varepsilon,\qquad v_2=-\varepsilon',
 \qquad \gamma(v_1)=1,\quad\gamma(v_2)=-1.
\]
On the local torus family over a linearizing disc \(\Delta_j\), the affine
map
\[
 (z,x)\longmapsto(g_jz,A_jx+v_j/m_j)
\tag{2.9}
\]
generates a free action of order \(m_j\).  The quotient
\(f_j:N_j\to D_j\) is a smooth complex threefold with
\[
 f_j^*(p_j)=m_jS_j,\qquad
 S_j=A_0^{(j)}/\langle x\mapsto A_jx+v_j/m_j\rangle.
\tag{2.10}
\]
Over the punctured disc it is biholomorphic to the original period family.
\end{proposition}

\begin{proof}
First prove freeness on the actual local family.  For a finite-order
integral map \(A\), the affine self-map \(x\mapsto Ax+b\) of
\((\Lambda\otimes\mathbb R)/\Lambda\) has a fixed point exactly when
\[
 \varphi(b)\in\mathbb Z
 \quad\text{for every }\varphi\in\operatorname{Hom}(\Lambda,\mathbb Z)
 \text{ with }\varphi\circ A=\varphi.
 \tag{2.9a}
\]
Necessity is immediate.  Conversely, the averaging projector
\(p=m^{-1}\sum_{r=0}^{m-1}A^r\) has image \(\ker(A-I)_{\mathbb R}\) and
kernel \(\operatorname{im}(A-I)_{\mathbb R}\); the dual lattice of
\(p(\Lambda)\) is exactly the invariant integral covector lattice.  Thus
(2.9a) puts \(pb\) in \(p(\Lambda)\), and
\(b\in(A-I)(\Lambda\otimes\mathbb R)+\Lambda\), which supplies a fixed
point.

The exact matrix calculation gives
\[
 V^{T_1^k}=\langle\gamma,\psi_1\rangle,
 \quad \psi_1=2u+w+3\delta \quad(k=1,2),
\]
and
\[
 V^{T_2^k}=\langle\gamma,\psi_2\rangle,
 \quad \psi_2=u+w+2\delta \quad(k=1,2,3).
 \tag{2.9b}
\]
For \(v_1=\varepsilon\) and \(v_2=-\varepsilon'\),
\[
 \gamma(v_1)=1,\quad\psi_1(v_1)=0,
 \qquad
 \gamma(v_2)=-1,\quad\psi_2(v_2)=0.
 \tag{2.9c}
\]
Because \(A_jv_j=v_j\), formula (2.8f) shows that the \(k\)-th power of
the logarithmic lift has translation \(kv_j/m_j\).  For \(j=1\), (2.9a)
fails for \(k=1,2\) already on \(\gamma\), since \(k/3\notin\mathbb Z\).
For \(j=2\), it fails for \(k=1,2,3\), since \(-k/4\notin\mathbb Z\).
Away from \(z_j\), the base map \(g_j^k\) has no fixed point in
\(\Delta_j\).  Hence every nonidentity power is fixed-point free on
\(J'_j\).  At \(k=m_j\), (2.8f) is translation by
\(v_j\in\Lambda\), hence the identity; the order is exactly \(m_j\)
because the map covers \(g_j\).

It follows that
\[
 N_j=J'_j/\langle\widetilde g_j^{\log}\rangle
\]
is a smooth complex threefold and the quotient map is finite \'etale of
degree \(m_j\).  The function \(s_j^{m_j}=t_j\circ\pi\) is invariant and
descends to \(f_j:N_j\to D_j\).  The projection
\(J'_j\to\Delta_j\) is proper: over a compact subset, the image of the
compact set obtained by letting \(x\) range over a fixed unit
parallelotope in \(\Lambda\otimes\mathbb R\) is the whole inverse image.
Together with finiteness of \(\Delta_j\to D_j\) and of the quotient, this
proves properness of \(f_j\).  Pulling the divisor of \(t_j\) back through
the etale quotient gives the divisor of \(s_j^{m_j}\), namely \(m_j\)
times the reduced central torus.  This proves (2.10).

It remains to identify the punctured filling with the original family and
to check that no logarithm branch has entered the answer.  On a simply
connected \(U\subset\Delta_j^*\), choose a branch \(\ell\) of
\(\log s_j\) and put
\[
 \sigma_j^\ell(z)=\frac{\ell(z)}{2\pi i}\Pi(z)v_j.
 \tag{2.9d}
\]
Changing \(\ell\) by \(2\pi in\) changes (2.9d) by the lattice vector
\(n\Pi(z)v_j\).  The classes therefore glue to a branch-independent
holomorphic section \(\sigma_j\) of \(J'_j|_{\Delta_j^*}\); translation by
that section defines a biholomorphism \(h_j\).

For the conjugacy, use \(\ell\) on \(U\) and on \(g_jU\) use the branch
\[
 \ell'(w)=\ell(g_j^{-1}w)-\frac{2\pi i}{m_j}.
\]
This is a logarithm of \(s_j(w)\) because
\(s_j\circ g_j=e^{-2\pi i/m_j}s_j\).  For \(z\in U\), equation (2.8e)
and \(A_jv_j=v_j\) give
\[
\begin{aligned}
 \widetilde g_j^{\log}h_j^{-1}[z,\zeta]
  =\Bigl[g_jz,R_{g_j}(z)\zeta
  -\Bigl(\frac{\ell(z)}{2\pi i}-\frac1{m_j}\Bigr)
    \Pi(g_jz)v_j\Bigr].
\end{aligned}
\]
But
\[
 \frac{\ell'(g_jz)}{2\pi i}
 =\frac{\ell(z)}{2\pi i}-\frac1{m_j},
\]
so applying \(h_j\) cancels the displayed translation and proves the
global identity
\[
 h_j\widetilde g_j^{\log}h_j^{-1}=\widetilde g_j^{0}.
 \tag{2.9e}
\]
Finally, by the stabilizer and translate conditions in (2.8b),
\[
 J|_{D_j^*}
 =\bigl(J'_j|_{\Delta_j^*}\bigr)/
   \langle\widetilde g_j^{0}\rangle.
\]
Thus (2.9e) descends to the asserted punctured-disc biholomorphism.
\end{proof}

\begin{lemma}[radial retraction and the finite-fibre stalk map]
\label{lem:finite-filling-stalk-map}
\Verified\ Under \cref{prop:candidate-finite-fillings}, the finite local
model \(N_j\) strongly deformation retracts onto its reduced central fibre
\(S_j\).  After identifying a nearby fibre with the marked real torus
\(A_s=(\Lambda\otimes\mathbb R)/\Lambda\) by the flat coordinates
\(\Xi\), the specialization map
\[
 H^r(S_j;\mathbb Z)\longrightarrow H^r(A_s;\mathbb Z)
\]
is the pullback \(\pi_j^*\) of the quotient covering
\[
 \pi_j:A_0^{(j)}\longrightarrow
 S_j=A_0^{(j)}/\langle x\mapsto A_jx+v_j/m_j\rangle .
\]
\end{lemma}

\begin{proof}
Use the linearizing coordinate to define the radial contraction
\(r_u:\Delta_j\to\Delta_j\) by
\[
 s_j(r_u(z))=(1-u)s_j(z),\qquad 0\le u\le1.
\]
Because \(s_j\circ g_j=\zeta_js_j\), this contraction commutes with
\(g_j\).  In the global real-analytic marking
\[
 \Xi(z,x+\Lambda)=[z,\Pi(z)x],
\]
put
\[
 H_u\bigl(\Xi(z,x+\Lambda)\bigr)
   =\Xi(r_u(z),x+\Lambda).
\]
Formula (2.8f) shows directly that \(H_u\) commutes with the logarithmic
lift: it keeps \(x\) fixed while both composites send it to
\(A_jx+v_j/m_j\).  Hence \(H_u\) descends to a strong deformation
retraction of \(N_j\) onto \(S_j\).

Fix \(t_j\ne0\) and one root \(s\) with \(s^{m_j}=t_j\).  Every orbit of
the logarithmic lift meets the fibre over \(s\) exactly once, so that fibre
maps isomorphically onto \(f_j^{-1}(t_j)\).  Following the descended radial
retraction sends the marked point \(x\in A_s\) to the orbit of the same
\(x\) in the central torus, namely \(\pi_j(x)\).  The induced map from
central-fibre cohomology to nearby-fibre cohomology is therefore exactly
\(\pi_j^*\), with no unrecorded conjugacy or index.
\end{proof}

\begin{theorem}[global analytic construction]
\label{thm:candidate-global-geometry}
\Verified\ Gluing the period family over \(B^\circ\) to
\(N_0,N_1,N_2\) by the preceding punctured-disc biholomorphisms produces a
connected compact Hausdorff complex threefold \(X\) and a proper surjective
holomorphic map
\[
 f:X\longrightarrow B\simeq\mathbb P^1
\]
with connected fibres.  The fibre over \(p_0\) is exactly the reduced
nonnormal surface \(W\) of \cref{thm:candidate-cusp-geometry}; those over
\(p_1,p_2\) have scheme multiplicities three and four.  The open embedding
\(N_0\hookrightarrow X\) identifies canonical bundles, so
\begin{equation}
 \omega_X|_{N_0}\simeq\omega_{N_0}\simeq\mathcal O_{N_0}.
\label{eq:candidate-global-canonical-on-cusp}
\end{equation}
\end{theorem}

\begin{proof}
Take the disjoint union of the four manifolds and identify the two copies of
each punctured-disc restriction by the displayed biholomorphism.  The special
discs are pairwise disjoint, so every nonsingleton equivalence class has two
points and meets each summand at most once.  The saturation of an open set is
the union of that set with its images under finitely many overlap
biholomorphisms, hence is open.  Each summand therefore embeds openly in the
quotient, and the only transition maps are biholomorphisms.  This gives a
complex threefold.

Points over distinct base values have disjoint disc preimages.  Points over
one base value lie in a common one of the four Hausdorff charts, so the
quotient is Hausdorff.  The local maps to the base agree on overlaps.  Each
is proper, for the period-family piece by
\cref{lem:middle-period-proper} and for the three fillings by their local
proper constructions; properness is local on a locally compact Hausdorff base, so the
global map is proper.  Since the base is compact, \(X=f^{-1}(B)\) is compact.
All four pieces are connected and every local filling meets the central
period-family piece, so \(X\) is connected.  The fibre statements follow
from \cref{thm:candidate-cusp-geometry,prop:candidate-finite-fillings}; their
fibres are connected, and surjectivity follows.
Because each summand embeds openly in the quotient, the restriction of the
intrinsic canonical bundle \(\omega_X\) to the open summand \(N_0\) is
\(\omega_{N_0}\).  Equation~(2.6a) proves
\cref{eq:candidate-global-canonical-on-cusp}.
\end{proof}

\paragraph{What has and has not been discharged.}
The formulas above reconstruct the period functions, both finite fillings,
the toric cusp quotient, and all three punctured-overlap biholomorphisms.
Thus \cref{thm:candidate-cusp-geometry,thm:candidate-global-geometry}
discharge existence, smoothness, properness,
the full central-fibre equation, and the normalization assertions formerly
labelled (H1)--(H3) in \cref{thm:candidate-conductor}.  This promotion is
local to the analytic construction: no collapse, Leray transgression, or
Mayer--Vietoris computation is used here.  Integral cohomology and sphere
recognition remain separate below.
